\title{Byte Latent Transformer: Patches Scale Better Than Tokens}

\begin{abstract}
We present the Byte Latent Transformer ({\textsc BLT}), a novel large language model architecture operating at the byte level that, for the first time, achieves performance comparable to traditional tokenization-based models at scale, while offering notable gains in inference speed and resilience. {\textsc BLT} transforms input bytes into adaptively sized segments, called patches, which act as the principal computational elements. The patch boundaries are determined by assessing the entropy of the subsequent byte, enabling the allocation of greater computational resources and model capacity in regions of higher data complexity. We perform the first flop-controlled scaling analysis of byte-level language models, evaluating architectures up to 8B parameters and training on 4T bytes. The findings verify that scaling language models directly from byte sequences—eschewing predefined vocabularies—is practical and effective. Both training and inference benefit from adaptively selecting extended patches where data exhibits high predictability, leading to improved efficiency and also enhancements in complex reasoning and coverage of rare phenomena. In summary, {\textsc BLT} achieves superior scaling at a constant inference budget compared to tokenization-based models, as it concurrently enlarges both patch length and model size.
\end{abstract}

\section{Introduction}
\label{section:intro}

We present the Byte Latent Transformer~(\textbf{BLT}), an innovative architecture that forgoes tokenization, instead directly processing unsegmented byte sequences. Remarkably, BLT attains parity in performance with conventional tokenization-dependent models at large scale, while also delivering notable gains in both computational efficiency and model robustness (\S\ref{sec:robustness}). In contrast, mainstream large language models (llms) are primarily trained in an end-to-end manner with the exception of an initial tokenization phase—a rule-based pre-processing operation that maps raw bytes to a predetermined vocabulary of tokens. This artificial segmentation impacts the way string data is compressed, contributing to drawbacks such as vulnerability to the specific domain or input type~\citep{dagangetting}, susceptibility to various forms of input corruption~(\S\ref{sec:robustness}), limited understanding of writing system structure~\citep{edman2024cute}, and inequities across languages~\citep{liang2023xlm,petrov2024language,limisiewicz2024myte}.

Tokenization has historically played a crucial role since training large language models on raw bytes incurs significant computational costs, largely due to the extended sequence lengths involved~\citep{xue2022byt5}. Earlier research has sought to alleviate this bottleneck through approaches such as more computationally efficient self-attention mechanisms~\citep{el2020characterbert, clark2022canine} or by leveraging architectures that entirely circumvent attention~\citep{wang2024mambabyte}~(\S\ref{section:related_work}). Nonetheless, these strategies yield their greatest benefits primarily for \textit{small models}. When considering scaling to larger models, it is the expansive feed-forward network layers—executed across every byte—that become the primary driver of computational expense, rather than the attention components.

To optimize compute usage, we introduce a flexible, trainable approach for assembling bytes into \textit{patches}~(\S\ref{section:patching}), along with a novel architecture capable of jointly processing both byte-level and patch-level data. Distinct from tokenization, {\textsc BLT} does not rely on a pre-defined patch vocabulary. Instead, any sequence of bytes can be transformed into a latent patch representation through compact, learnable encoder and decoder components. Our results demonstrate that this yields \textit{greater} computational efficiency compared to traditional tokenization-driven models.

Tokenization-based llms assign an equal amount of computational resources to each token, leading to a compromise between computational efficiency and model performance. This uniform allocation stems from token compression strategies that may not accurately reflect the actual prediction complexity. Our approach centers on the principle that computational effort should be distributed adaptively, depending on necessity. For instance, it is generally unnecessary to employ a large transformer for predicting the final characters in words, as these typically involve straightforward, low-entropy choices, unlike the selection of the initial token in a sentence, which presents a greater challenge. This principle underlies the design of {\textsc BLT}~(\S\ref{section:architecture}), which utilizes three transformer modules: two lightweight byte-level \textit{local models} and a more substantial global \textit{latent transformer}~(\autoref{fig:arch_high_level}). 

To facilitate adaptive allocation of computation by grouping bytes into patches, {\textsc BLT} segments input data according to the entropy of predicting the subsequent byte, yielding contextualized clusters of bytes with roughly consistent informational content.

We conduct the inaugural flop-controlled scaling analysis of byte-level models up to 8B parameters and 4T bytes of training data, demonstrating that large-scale end-to-end training directly from bytes—bypassing fixed-vocabulary tokenization—is viable. In comprehensive experiments, {\textsc BLT} achieves parity with Llama 3 on training flop-controlled benchmarks\footnote{Model computational expense is assessed via the total Floating Point OPerations (flops) executed.} while reducing inference flops by as much as 50\% (\S\ref{section:scaling}). 

Moreover, our results indicate that operating on raw bytes substantially enhances performance on infrequent data distributions. Compared to models reliant on tokenizers, {\textsc BLT} exhibits greater resilience to noisy input and superior understanding of characters, as evidenced by evaluations in orthography, phonology, and low-resource machine translation tasks~(\S\ref{sec:robustness}).

Additionally, {\textsc BLT} facilitates the concurrent enlargement of both model capacity and patch size without increasing the inference flop requirement. Utilization of larger patch sizes generally reduces computation per inference, permitting reallocation of compute to scale up the global latent transformer, since it executes less frequently. Through inference-flop controlled scaling investigations (\autoref{fig:fixed_inference_scaling}), we find markedly improved scaling characteristics compared to conventional tokenization-based approaches.

To conclude, the key contributions of this work are as follows: 1) We present {\textsc BLT}, a byte-level latent large language model architecture that dynamically adjusts computational resources, thereby enhancing flop efficiency; 2) We demonstrate that our approach achieves parity with Llama 3 in terms of training flops up to an 8B model size, while also allowing for up to a 50\% reduction in flops by making minor compromises in evaluation outcomes; 3) {\textsc BLT} architectures enable a novel approach to scaling llms, as they permit increases in model capacity while keeping inference costs constant; and 4) We provide evidence that {\textsc BLT} models exhibit greater resilience to noisy inputs and exhibit a nuanced understanding of sub-word information in data, outperforming token-based llms in this regard. The full training and inference implementation for {\textsc BLT} is openly available at~\url{https://github.com/facebookresearch/blt}.

\section{Patching: From Individual Bytes to Groups of Bytes}
\label{section:patching}

Dividing byte sequences into distinct \textit{patches} enables {\textsc BLT} to flexibly distribute computational resources in response to contextual demands. Figure~\ref{fig:patching_types} illustrates various strategies for segmenting bytes into patches. More formally, a patching function $f_p$ partitions an input byte sequence $\pmb{x}=\{x_i,|i=1,\ldots n\}$ of length $n$ into a series of $m < n$ patches $\pmb{p}=\{p_j|j=1,\ldots,m\}$, where each byte $x_i$ is assigned a value in \{0,1\}, with a 1 marking the initiation of a new patch. The overall computation required for both token-based and patch-based architectures is primarily a function of the total number of steps performed by the principal Transformer component. In the case of {\textsc BLT}, this equates to the count of patches generated to represent the input data under the particular patching function in use. As such, the typical patch length—or \textit{patch size}—serves as the principal metric influencing the computational cost associated with both training and inference for any specific patching method~(\S\ref{section:flops}).

We subsequently present three distinct patching schemes: dividing the input into patches of a predetermined byte count~(\S\ref{section:static-patch}), delineating patches using whitespace boundaries~(\S\ref{section:white-patch}), and assigning patches adaptively based on byte-level entropy computed via a lightweight language model~(\S\ref{section:dyn-patch}). Lastly, we elaborate on incremental patching strategies and clarify distinctions between patching and tokenization~(\S\ref{section:bpe-patch}).

\subsection{Strided Patching Every K Bytes}
\label{section:static-patch}

A common approach to grouping bytes involves dividing them into uniformly sized patches of $k$ bytes, as seen in MegaByte~\citep{yu2023megabyte}. Employing a fixed stride in this manner is straightforward to implement during both training and inference phases, enables simple tuning of the average patch size, and consequently offers a clear method to regulate computational cost. Nonetheless, this fixed patching method introduces notable limitations. Firstly, computational resources are not adaptively focused where they might be most effective: a transformer step $j$ may be inefficiently used, for example, by being assigned to predict only whitespace in code, while failing to allocate adequate capacity for highly informative byte regions such as mathematical content. Secondly, such fixed patching results in inconsistent and context-insensitive treatment of similar byte sequences; the same word, for instance, may be segmented in divergent ways.

\subsection{Space Patching}
\label{section:white-patch}

\citet{slagle2024spacebyte} introduces a straightforward yet powerful modification to strided patching by generating new patches at each space-like byte\footnote{
    Space-like bytes refer to bytes that are neither Latin characters, digits, nor UTF-8 continuation bytes. Furthermore, every patch is required to include at least one non space-like byte.
}—locations which often represent the boundaries between linguistic units across various languages. Under the Space patching approach, every word is afforded an additional latent transformer computation step (incurring extra flops), enabling consistent patching of words within and across sequences and ensuring computational resources are directed at difficult predictions, which frequently arise after space characters. For instance, anticipating the initial byte of a response in ``Who composed the Magic Flute?\uline{\hspace{.6em}}'' poses a substantially greater challenge than determining the subsequent bytes in ``Mozart'' after ``M.'' This is because the first letter dramatically narrows the possible completions, making subsequent predictions much more straightforward. Nonetheless, the space patching strategy exhibits limitations: it is not universally applicable across all languages or data types, and it lacks the ability to flexibly adjust patch size. In what follows, we propose a novel patching technique that builds on the idea that word-initial bytes generally present the greatest difficulty in prediction, while also inherently allowing for dynamic control over patch sizes.

\subsection{Entropy Patching: Using Next-Byte Entropies from a Small Byte LM}
\label{section:dyn-patch}

Instead of utilizing rule-based heuristics like whitespace to determine boundaries, we employ a data-driven strategy for detecting points of high uncertainty in next-byte predictions. Specifically, we propose \textit{entropy patching}, a method that leverages entropy estimates to establish patch boundaries.

To this end, we train a compact auto-regressive language model at the byte level using the {\textsc BLT} training set. We then evaluate the entropy for each potential next byte based on the language model's predictive distribution $p_{e}$ over the byte-level vocabulary $\mathcal{V}$:

\begin{equation}
H(x_i) = \sum_{v \in \mathcal{V}} p_{e}(x_i=v|\pmb{x}_{<i}) \log p_{e}(x_i=v|\pmb{x}_{<i})
\end{equation}

We explore two distinct strategies for determining patch boundaries utilizing entropy values $H(x_i)$. The first approach selects points that surpass a predefined global entropy threshold, as shown in~\autoref{fig:patching}. The second approach targets points where entropy is notably higher in comparison to the preceding value. This latter method may also be viewed as locating points that disrupt an otherwise approximately monotonically decreasing pattern of entropy within the patch.

\begin{align*}
    \text{Global Constraint}\hspace{1cm}H(x_t)& > \theta_g\\
    \text{Approx. Monotonic Constraint}\hspace{1cm}H(x_t)& - H(x_{t-1}) > \theta_r
\end{align*}

Patch boundaries are determined as part of an efficient preprocessing phase carried out when data is being loaded. This contrasts with~\citet{nawrot-etal-2023-efficient}, in which a classifier is trained specifically to identify entropy-based patch boundaries. Within our experiments~(\S\ref{section:experiments}), we conduct a comparison between these two techniques for discriminating low entropy from high entropy bytes.

\subsection{The Byte-Pair Encoding (BPE) Tokenizer and Incremental Patching}
\label{section:incremental-patching}
\label{section:bpe-patch}

Contemporary llms, such as our baseline model Llama 3, typically employ subword tokenization methods like bpe~\citep{gage1994new, sennrich-etal-2016-neural}. In this context, we define ``tokens'' as byte sequences selected from a \textit{finite} vocabulary established before the commencement of training, distinguishing them from ``patches,'' which are flexibly constructed groupings of bytes lacking a predetermined vocabulary. A fundamental distinction is that when models work with tokens, they are not exposed to the original byte-level features directly.

An important advancement introduced by {\textsc BLT} relative to tokenization-based approaches lies in its reconceptualization of the balance between vocabulary size and computational demands. In conventional language models, expanding the vocabulary leads to longer tokens on average, resulting in reduced sequence length during inference but necessitating a larger output dimension in the model’s final projection layer. This inherent trade off significantly limits the flexibility of tokenization-based strategies to substantially alter token lengths or reduce inference expenses. As an illustration, Llama 3 increases the mean token length from 3.7 to 4.4 bytes, but this comes at the expense of enlarging the embedding table by a factor of four compared to Llama 2.

During generation, {\textsc BLT} must determine whether the present location within the byte sequence aligns with a patch boundary, as this dictates whether additional computation through the Latent Transformer is necessary. Importantly, this assessment must be made without knowledge of forthcoming elements in the sequence, meaning that segmentation into patches cannot rely on access to future bytes. In formal terms, we define an incremental patching scheme $f_p$ as one for which the following property holds:
$$f_p(\pmb{x}_{<i}) = f_p(\pmb{x})_{<i}$$
The commonly-used bpe approach fails to be an incremental patching scheme, since the way a prefix is tokenized can depend on subsequent bytes in the sequence; consequently, it does not fulfill the condition stated above\footnote{Introducing a dedicated delimiter to signify patch boundaries may render bpe incremental, but it comes at the cost of lengthening the byte sequence.}.

\section{{\textsc BLT} Architecture}
\label{section:architecture}
The {\textsc BLT} framework integrates a large-scale, global autoregressive language model that processes representations at the patch level. Additionally, it leverages two more compact, local modules that respectively handle the encoding of byte sequences into patch representations and the decoding of patch representations back to bytes, as depicted in~(\autoref{fig:arch_high_level}).

\subsection{Latent Global Transformer Model}

\textit{The Latent Global Transformer} refers to an autoregressive transformer architecture, denoted $\mathcal{G}$ and comprising $l_{\mathcal{G}}$ layers, that transforms a series of latent input patch embeddings $p_j$ into corresponding output patch representations $o_j$. In this manuscript, the index $j$ is reserved for patches, while $i$ indexes bytes. The global transformer incorporates a block-causal attention mask~\citep{dubey2024llama}, ensuring that attention operations are confined to the current patch and preceding patches within the same document. As the principal consumer of floating point operations (flops) during both pre-training and inference, this model’s computational demands can be managed by strategically selecting when it is utilized. Consequently, the computational cost for processing different regions of the input and output can be dynamically adjusted according to their respective complexities.

\subsection{Local Encoder}
\label{section:local}

\textit{The Local Encoder Model}, represented as $\mathcal{E}$, utilizes a compact transformer architecture with significantly fewer layers, $l_{\mathcal{E}} << l_{\mathcal{G}}$, in comparison to the global model. Its core function is to rapidly encode input byte sequences $b_i$ into robust patch-level representations $p_j$. One fundamental modification to the typical transformer structure is the insertion of a cross-attention module following every transformer block, enabling the aggregation of byte-level features into patch-level representations~(\autoref{fig:crossattn}).

The initial step involves mapping each byte in the input sequence, $b_i$, into an embedded vector $x_i$ through a learnable matrix with shape $\mathbb{R}^{256 \times h_{\mathcal{E}}}$. Optionally, these initial embeddings can be enhanced with hash-based embeddings as described in~(\S\ref{sec:hash-emb}). Subsequently, alternating stacks of transformer and cross-attention blocks~(\S\ref{sec:enc-cross-attn}) refine these vectors, ultimately yielding patch representations, $p_i$, to be consumed by the global transformer $\mathcal{G}$.

Within each transformer block, \textit{local block causal} attention masking is adopted: a given byte is permitted to attend over a limited window of $w_{\mathcal{E}}$ preceding bytes, potentially spanning dynamic patch borders but never crossing document boundaries. The following sections further elaborate on the specifics of both the embedding scheme and the cross-attention architecture.

\subsubsection{Encoder Hash n-gram Embeddings}
\label{sec:ngram-emb}
\label{sec:hash-emb}
An essential factor in generating resilient and informative representations at every step $i$ is integrating context from previous bytes. In {\textsc BLT}, this is accomplished by capturing the representation of the byte $b_i$ on its own as well as within the context of a surrounding byte n-gram. Specifically, at each step $i$, we begin by forming byte-grams

\begin{align}
    g_{i,n}=\{b_{i-n + 1},\ldots, b_i\}
\end{align}

for each byte position $i$ and $n$ from three to eight.\footnote{
We omit byte-grams of size $n$ or more when $i<n$. }

Next, we present hash $n$-gram embeddings, which assign each byte $n$-gram, for $n \in\{3,4,5,6,7,8\}$, to an entry in a fixed-size embedding table $E_{n}^{hash}$ using a hash function~\citep{bai2010learning}. The embedding derived from this process is combined with the original byte embedding, followed by normalization, before being fed into the local encoder model. The enriched embedding is computed

\begin{align}
e_i &= x_i + \sum_{n = 3,...,8}E_{n}^{hash}(\text{Hash}(g_{i,n})) \\
\text{where, Hash}(g_{i,n}) &= \text{RollPolyHash}(g_{i,n}) \% |E_{n}^{hash}|
\end{align}

To account for the varying number of $n$-gram sizes, we divide $e_i$ by the total count of $n$-gram sizes incremented by one, employing RollPolyHash as outlined in Appendix~\ref{appendix:rollpolyhash}. In Section~\ref{section:ablations}, we systematically investigate how altering the value of $n$ and the size of the embedding table for $n$-gram hash embeddings influences flop-controlled scaling law behavior. Beyond hash-based $n$-gram embeddings, we additionally assessed frequency-based $n$-gram embeddings; further information regarding these experiments is presented in Appendix~\ref{appendix:ngrams}.

\subsubsection{Encoder Multi-Headed Cross-Attention}
\label{sec:enc-cross-attn}
Our approach largely mirrors the input cross-attention mechanism outlined in the Perceiver architecture~\citep{jaegle2021perceiver}, except that, in our method, the latent vectors align with representations of variable patches, rather than utilizing a predetermined set of latents~(\autoref{fig:crossattn}). Furthermore, these representations interact solely with the byte sequences belonging to their specific patch. For each patch $p_j$, we establish a query representation by applying a pooling operation over the associated byte embeddings, which is then projected through a linear mapping, $\mathcal{E}_{C} \in \mathbb{R}^{h_{\mathcal{E}} \times (h_{\mathcal{E}} \times U_{\mathcal{E}})}$, with $U_{\mathcal{E}}$ denoting the total encoder cross-attention heads. Concretely, letting $f_{\text{bytes}}(p_j)$ indicate the bytes for patch $p_j$, we compute

\begin{align}
P_{0,j} &= \mathcal{E}_{C}(f_\text{bytes}((p_j)), f~ \text{is a pooling function} \\
P_l &= P_{l-1} + W_o\left(\text{softmax}\left(\frac{QK^T}{\sqrt{d_k}}\right)V\right) \\
\text{where } Q_j &= W_q(P_{l-1,j}), K_i = W_k(h_{l-1,i}), V_i = W_v(h_{l-1,i}) \\
h_l &= \text{Encoder-Transformer-Layer}_l(h_{l-1})
\end{align}

Here, $P \in \mathbb{R}^{n_p \times h_{\mathcal{G}}}$ denotes the collection of $n_p$ patch representations that serve as input to the global model; these are constructed by aggregating the byte embeddings $e_i$ associated with each patch $p_j$. The matrices $W_q$, $W_k$, $W_v$, and $W_o$ correspond to the transformations applied to generate queries, keys, values, and output, respectively, with the keys and values being computed as linear projections of the byte representations $h_i$ from the preceding layer (or from $e_i$ if in the initial layer). A patch-specific masking approach is employed such that for each query $Q_j$, attention is restricted solely to those keys and values linked to the bytes in patch $j$. Since multi-headed attention is performed over $Q, K,$ and $V$, and patch representations generally possess higher dimensionality ($h_{\mathcal{G}}$) compared to $h_{\mathcal{E}}$, during cross-attention we retain $P_l$ as a set of attention heads each of size $h_{\mathcal{E}}$, subsequently concatenating these to obtain the full $h_{\mathcal{G}}$-dimensional representation. Furthermore, we implement pre-LayerNorm on the queries, keys, and values; this cross-attention mechanism does not utilize positional embeddings. A residual connection wraps the entire cross-attention module.

\subsection{Local Decoder} 

Analogous to the structure of the local encoder, the local decoder $\mathcal{D}$ utilizes a lightweight transformer framework comprising $l_{\mathcal{D}}$, where $l_{\mathcal{D}} << l_{\mathcal{G}}$, transformer layers. Its function is to reconstruct a sequence of raw bytes, $y_i$, from the sequence of global patch embeddings $o_j$. During this process, the local decoder models each byte conditioned on the sequence of previously generated bytes, relying on the representations generated by the local encoder for the input byte sequence. The architecture consists of $l_{\mathcal{D}}$ layers that alternate between cross-attention and transformer modules. Notably, the cross-attention operation is applied first within each layer, transforming global patch representations into byte-level representations, which are subsequently refined by the decoder's transformer layer acting directly on the byte sequence.

\subsubsection{Decoder Multi-headed Cross-Attention} 

Within the decoder cross-attention mechanism, the assignments of queries and keys/values are swapped: the byte-representations function as the queries, while the patch representations serve as keys and values. At the start of cross-attention, byte-representations are drawn from the byte embeddings output by the final encoder layer, denoted $h_{l_{\mathcal{E}}}$. For each subsequent layer $l$, the byte-representations $d_{l,i}$ are updated using the following process:

\begin{align}
D_0 &= h_{l_{\mathcal{E}}} \\
B_l &= D_{l-1} + W_o\left(\text{softmax}\left(\frac{QK^T}{\sqrt{d_k}}\right)V\right), \\
  &\text{where } Q_i = W_q(d_{l-1,i}), K_i = W_k(\mathcal{D}_C(o_j)), V_i = W_v(\mathcal{D}_C(o_j)) \\
  D_l &= \text{Decoder-Transformer-layer}_l(B_l)
\end{align}

In this context, $W_k$ and $W_v$ once again refer to the key and value projection matrices, which implement a linear transformation and a split operation $\mathcal{D}_C$ on the output patch representations $o_j$ from the global model. The query projection matrix $W_q$ acts upon the byte-level representations $d_{l-1}$ produced by the prior decoder transformer layer (or $h_{l_{\mathcal{E}}}$ if it is the initial layer), and $W_o$ designates the output projection matrix. Together, these ensure that $B \in \mathbb{R}^{h_{\mathcal{D}} \times n_b}$, where $n_b$ denotes the quantity of output bytes. The subsequent decoder states $D_l$ result from applying a decoder transformer layer to the cross-attention output $B$. We adopt the same strategy as in the local encoder cross-attention by leveraging multi-head attention, incorporating pre LayerNorm, omitting positional embeddings, and implementing a residual connection around the cross-attention component.

\section{Experimental Setup}
\label{section:experiments}
We systematically construct a suite of controlled experiments to rigorously evaluate {\textsc BLT} against models that use traditional tokenization, ensuring that no inadvertent advantage arises for {\textsc BLT} due to potential access to longer sequence contexts.

\subsection{Pre-training Datasets} In our study, all model sizes are initially trained on two primary datasets: (1) The Llama 2 dataset~\citep{touvron2023llama}, containing 2 trillion tokens sourced from a diverse array of publicly available materials, followed by comprehensive cleaning and filtering to ensure high data quality; and (2) {\textsc BLT}-1T, a newly assembled corpus totaling 1 trillion tokens drawn from numerous public data repositories, which also incorporates a subset of the pre-training data provided by Datacomp-LM~\citep{li2024datacomp}. For our scaling law analyses—focusing on optimal token counts as identified by~\cite{dubey2024llama}—we rely on the Llama 2 dataset to assess ideal architectural parameters for {\textsc BLT}. In contrast, {\textsc BLT}-1T is reserved for conducting a full pre-training cycle, facilitating direct comparisons with Llama 3 on a range of downstream evaluation tasks. It is important to note that neither dataset contains any information harvested from Meta’s products or services. Additionally, to establish baseline metrics in our tokenizer-based model experiments, we make use of the Llama 3 tokenizer, which features a vocabulary of 128K tokens and has empirically demonstrated improved baseline results over the Llama 2 tokenizer in our tests.

\subsection{Entropy Model} For our experiments, the entropy model comprises a byte-level language model, which is trained using the same training dataset as the full {\textsc BLT} model. By default, we employ a transformer architecture featuring 100 million parameters, a depth of 14 layers, and a hidden size of 512, utilizing a sliding window attention mechanism across 512 bytes. All other hyperparameters mirror those set for our local and global transformer variants. As outlined in \autoref{section:ablations}, we explored variations in model size, receptive field length, and overall architecture. Notably, in cases where the receptive field is sufficiently limited, it becomes possible to encode the learned entropy model as a compact lookup table.

\subsection{Entropy Threshold and Equalizing Context Length} For entropy-based patching models, we determine a patching threshold corresponding to a specified average \textit{patch size} within the pretraining data blend. In {\textsc BLT}, in contrast to fixed tokenization, the \textit{patch size} is not constrained and can be freely selected, which directly affects the model's effective context size. To ensure a fair comparison and prevent larger patch sizes from resulting in disproportionate context windows, we keep the total expected byte count per batch invariant. Consequently, for models that employ larger patch sizes, we proportionally decrease their sequence lengths. Specifically, for Llama 2 training data, we adopt an 8k byte context, whereas for the {\textsc BLT}-1T dataset, the context window is set to 16k bytes on average, all while the average batch size is fixed at 16M bytes.

Although the mean batch size remains fixed, dynamic patching approaches produce varying proportions of bytes per patch during batch data loading. To enhance efficiency, our implementation of {\textsc BLT} training groups patches into batches such that padding is not needed within the latent transformer, which is computationally intensive. This method maintains a uniform patch count across all batches. During the training process, we pad or, if necessary, truncate the byte sequences to 12k bytes for Llama 2 and to 24k bytes for the {\textsc BLT}-1T datasets, thereby preventing memory overload that could occur when sequences contain exceptionally large patches.

\subsection{Entropy Model Context}
\label{section:entropy-context}
Our experiments indicate that entropy patching leads to increasingly larger patches when applied to structured material, such as in multiple choice tasks (\autoref{fig:mmlu_patching}), where content tends to be highly repetitive. This pattern is attributable to reduced entropy for repeated segments within the entropy model context. Consequently, for the main {\textsc BLT}-Entropy experiment with a patch size of 4.5, we refresh the entropy context with new line delimiters and adopt the approximate monotonicity constraint, which demonstrates greater resistance to "entropy drift" resulting from variations in context length. This adjustment only influences our calculation of entropies; the approach for determining the entropy threshold remains unchanged.

\subsection{FLOPs Estimation} 
\label{section:flops}

Our methodology for calculating transformer FLOPs is primarily based on the formulas provided in Chinchilla~\citep{hoffmann2022training}, incorporating FLOPs contributions from feed-forward components, qkvo projections within the self-attention block, as well as attention score calculations and output projections. However, our approach diverges in that we model the input embedding layer as a computationally inexpensive lookup table rather than a dense matrix multiplication, resulting in zero FLOPs attributed to this step. Consistent with prior literature, we approximate that the backward propagation phase entails double the FLOPs of the forward propagation.

In order to determine the number of flops \textit{per byte} for {\textsc BLT} architectures, we sum the computational costs for the local encoder transformer, the global latent transformer, and the local decoder transformer, as well as the contributions from cross attention mechanisms present within both the encoder and decoder. This can be mathematically formulated as follows:

\begin{align}
    \text{FL}_{\text{{\textsc BLT}}} &= \text{Transf. FL}(h_{\mathcal{G}}, l_{\mathcal{G}}, m=n_{ctx}/n_p,V=0) / n_p\\
   &+\text{Transf. FL}(h_{\mathcal{E}}, l_{\mathcal{E}}, m=w_{\mathcal{E}}, V=0) \\
   &+ \text{Transf. FL}(h_{\mathcal{D}}, l_{\mathcal{D}}, m=w_{\mathcal{D}}, V=256) \\ 
    &+ \text{Cross Attn. FL}(h_{\mathcal{E}}, l_{\mathcal{E}}, m=n_p, r=n_p/k) \times k/n_p \\ 
   &+ \text{Cross Attn. FL}(h_{\mathcal{D}}, l_{\mathcal{D}}, m=k, r=k/n_p) 
\end{align}

In this context, $n_{ctx}$ represents the byte-level sequence length, while $n_p$ denotes the size of each patch. The parameter $r$ specifies the proportion of queries to key/values, and $k$ signifies the ratio between patch and byte dimensions; equivalently, it indicates how many local model segments are combined to construct the comprehensive global representation ($k=2$ as illustrated in \autoref{fig:crossattn}). $V$ indicates the vocabulary size pertinent to the output projection, utilized exclusively by the local decoder. The attention module adapts its context length, $m$, based on whether it operates over byte sequences or patch sequences. Correspondingly, we adjust the calculation of attention-related flops for each part. Detailed formulas for calculating the flops associated with both the Transformer and Cross-Attention modules can be found in Appendix~\ref{section:flops-appendix}.

\subsection{Bits-Per-Byte Estimation} Perplexity is meaningful only when used with a specific tokenizer, as it quantifies the unpredictability of individual tokens. In order to enable fair evaluation between models operating at the byte level and those at the token level, as established in prior studies~\citep{xue2022byt5, yu2023megabyte, wang2024mambabyte}, we utilize Bits-Per-Byte (BPB) as our metric. BPB provides a tokenizer-agnostic alternative to perplexity. Concretely:

\begin{align}
    \text{BPB}(x)&= \frac{\mathcal{L}_{CE}(\pmb{x})}{\ln(2)\cdot n_{\text{bytes}}}
\end{align}

where the uncertainty over the data $\pmb{x}$ as measured by the sum of the cross-entropy loss is normalized by the total number of bytes in $\pmb{x}$ and a constant. 

\subsection{Transformer Architecture Hyperparameters} 

Across all transformer blocks in {\textsc BLT}, encompassing both the local and global configurations, we adopt an architecture modeled closely on Llama~3~\citep{dubey2024llama}. Specifically, we implement the SwiGLU activation~\citep{swiglu} within the feed-forward modules, employ rotary positional embeddings (RoPE)~\citep{RoPe} parameterized with $\theta=500000$~\citep{xiong2024effective}—applying them exclusively to self-attention mechanisms—and utilize RMSNorm~\citep{RMSNorm} for normalization at each layer. For self-attention layers governed by fixed-standard attention schemes, such as \textit{block causal} or \textit{fixed-window block causal}, we use Flash attention~\citep{dao2022flashattention}, specifying a window of 512 tokens for fixed-width configurations. When handling cross-attention, where attention masks are adaptively determined based on patch information, we leverage Flex Attention\footnote{\url{https://pytorch.org/blog/flexattention}}. This approach yields fused kernel operations, which substantially accelerate training procedures.

\subsection{{\textsc BLT}-Specific Hyperparameters} In order to assess the performance of {\textsc BLT} models, we undertake experiments in two key areas: analyzing scaling behavior and evaluating on downstream tasks. For this purpose, we examine models of various sizes—400M, 1B, 2B, 4B, and 8B parameters. Detailed hyperparameters concerning model architecture can be found in Appendix~Table~\ref{tab:arch}. 
Within the local encoder, we apply max-pooling to set the initial queries for the first cross-attention layer. Across all {\textsc BLT} variants, a hash table comprising $500,000$ entries with a single hash function and n-gram lengths spanning 3 to 8 is employed. The learning rate is uniformly set to $4e-4$ across all scales. This decision to align the learning rates for token and {\textsc BLT} models is based on a hyperparameter sweep at the 400M and 1B parameter levels, exploring learning rates between $1e-3$ and $1e-4$ and revealing $4e-4$ to be consistently optimal. Batch sizes during scaling trend experiments on Llama-2 datasets are chosen as per recommendations in~\cite{dubey2024llama}, or are adjusted according to equivalent data volume in bytes. For the optimization procedure, we utilize the AdamW optimizer~\citep{adamw} with $\beta_1 = 0.9$, $\beta_2 = 0.95$, and $\epsilon=10^{-8}$. A learning rate schedule incorporating a linear warm-up over the initial 2000 steps and followed by cosine decay to zero is implemented. Additionally, a weight decay of 0.1 is enforced, and the global gradient norm is clipped to a maximum of 1.0.

\section{Scaling Trends}
\label{section:scaling}

We offer a comprehensive analysis of scaling behaviors in byte-level models to guide the continued scaling of {\textsc BLT} models. Our scaling investigation seeks to overcome several shortcomings found in earlier work on byte-level modeling by: (a) analyzing scaling patterns under compute-optimal training regimes, (b) training comparable 8B parameter models on substantial datasets (reaching up to 1T tokens or 4T bytes) and evaluating their downstream task performance, and (c) examining scaling properties while holding inference cost constant. Further discussion on unique benefits of modeling byte-sequences is provided in a subsequent section.

\subsection{Parameter Matched Compute Optimal Scaling Trends}
\label{section:parameter-matched-scaling}
Leveraging the Llama 2 dataset, we conduct training runs for both \textit{compute-optimal} bpe and {\textsc BLT} models at four different scales, with parameter counts spanning from 1B to 8B. We assess language modeling outcomes by plotting training flops against performance on a representative selection of the training data mixture. In alignment with the findings from Llama~3~\citep{dubey2024llama}, the bpe models are configured according to the empirically optimal proportion between parameter size and dataset volume. This \textit{compute-optimal} configuration is grounded in theoretical work that identifies the optimal model and data ratio for maximizing in-domain performance within a fixed computational budget~\citep{hoffmann2022training}, establishing a strong reference point for comparison. For every bpe model considered, a {\textsc BLT} model of identical scale and transformer architecture is trained on the same data, with a Latent Transformer whose parameters and structural properties precisely match its bpe counterpart.

As shown in~\autoref{fig:scaling-compute-opt} (right), {\textsc BLT} models consistently equal or surpass the performance of their bpe equivalents, and this pattern persists across increases in model size and flops. To our knowledge, {\textsc BLT} represents the first Transformer architecture operating at the byte level to demonstrate comparable scaling properties to BPE-based models when trained under compute-optimal conditions. This observation lends support to our hypothesis that the ideal ratio of parameters to training compute for bpe is also applicable to {\textsc BLT}, or at least does not differ substantially.

Both advancements in model architecture and the adoption of dynamic patching are essential to achieve scalability comparable to bpe trends. In ~\autoref{fig:scaling-compute-opt} (left), we present a comparison between models utilizing space-patching strategies and Llama 3. For this analysis, we approximate SpaceByte \citep{slagle2024spacebyte} by employing {\textsc BLT} space-patching, omitting the use of n-gram embeddings and cross-attention. While SpaceByte demonstrates improvements over Megabyte, it does not reach the performance level exhibited by Llama 3. \autoref{fig:scaling-compute-opt} (right) further visualizes how enhancements in model architecture in conjunction with dynamic patching contribute to overall performance. At scale, {\textsc BLT} models achieve results that are competitive with leading tokenizer-based systems such as Llama 3.

Additionally, our findings indicate that the specific tokenizer employed in training tokenizer-based models significantly impacts performance: models trained with the Llama-3 tokenizer consistently surpass those utilizing the Llama-2 tokenizer when evaluated on identical training datasets.

Ultimately, our {\textsc BLT} model exhibits a performance trajectory between that of Llama 2 and Llama 3 when evaluated with substantially larger patch sizes. While the average token produced by the bpe tokenizers in Llama 2 and 3 is 3.7 and 4.4 bytes respectively, {\textsc BLT} is able to sustain comparable scaling behavior using average patch sizes of 6 or even 8 bytes. Since inference flops decrease as average patch size increases, adopting an 8-byte patch size would nearly halve the inference flop requirements. Additionally, models employing larger patch sizes tend to demonstrate improved performance as both model and dataset size are scaled up. For example, while {\textsc BLT} with an 8-byte patch size begins at a clear disadvantage relative to bpe Llama 2 at the 1B parameter scale, it surpasses the bpe approach by the time model size reaches 7B. This observation points to the potential of larger patch sizes to yield superior results at greater model scales, suggesting the possibility of success with even more sizable patches as model capacity and available training computation increase.

\subsection{Beyond Compute Optimal Task Evaluations}
\label{section:evals}
To further investigate scaling behavior, we extend training of an 8B {\textsc BLT} model past the standard compute-optimal regime using the {\textsc BLT}-1T dataset, which comprises a larger and higher quality corpus. The model’s capabilities are then assessed on a range of widely-used classification and generative benchmarks. For this evaluation, tasks are chosen to reflect diverse abilities in commonsense reasoning, factual knowledge, and code synthesis:
\paragraph{Classification tasks} assessed are ARC-Easy (zero-shot)~\citep{Eval_ARC}, Arc-Challenge (zero-shot)~\citep{Eval_ARC}, HellaSwag (zero-shot)~\citep{Eval_hellaswag}, PIQA (zero-shot)~\citep{Eval_piqa}, and MMLU (five-shot)~\citep{hendrycks2020measuring}. We utilize prompt-based likelihood computation across candidate responses, reporting mean classification accuracy.
\paragraph{Coding related generation tasks:}  Performance is measured with pass@1 accuracy on MBPP (three-shot)~\citep{Eval_MBPP} and HumanEval (zero-shot)~\citep{Eval_HumanEval}, allowing us to evaluate large language models' proficiency in synthesizing Python programs.

In \autoref{tab:evals}, a comparison is presented among three models trained on the {\textsc BLT}-1T dataset: a bpe model using the Llama 3 tokenizer,\footnote{The Llama 3 tokenizer, with its 128k vocabulary, is selected due to its superior performance compared to the 32k vocabulary of Llama 2.} and two {\textsc BLT} model variants. The first employs a space-patching approach ({\textsc BLT}-Space), while the second applies an entropy-driven patching method ({\textsc BLT}-Entropy), implementing an approximate monotonicity constraint and resetting the entropy model's context at new lines, as explained in~\autoref{section:entropy-context}. All models receive equivalent computational resources in terms of flop budget for training. For the {\textsc BLT}-Entropy model, an additional modification at inference is made: the entropy threshold is lowered from 0.6 to 0.1, which is observed to enhance performance on evaluation tasks, though it requires additional inference steps.

The {\textsc BLT}-Entropy model achieves higher performance than the Llama 3 model on four of the seven evaluated tasks, despite both models being trained with an equal quantity of bytes. This enhancement can be attributed to two main factors: (1) more efficient utilization of training computation through dynamic patching, and (2) the explicit modeling of information at the byte level rather than at the token level.

Conversely, {\textsc BLT}-Space performs worse than the Llama 3 tokenizer on all but a single task, although its use of a larger average patch size of 6 bytes results in a notable decrease in inference flops. In contrast, models based on bpe and entropy-patching exhibit similar average patch sizes, approximately 4.5 bytes, when evaluated on the mixed training dataset. Given an identical training budget, the model utilizing the larger patch size processes 30\% more data than the other two approaches, which may drive {\textsc BLT} further from the compute-optimal regime.

\subsection{Patches Scale Better Than Tokens} {\textsc BLT} architectures enable concurrent scaling of both model capacity and patch dimensions, all while preserving a constant budget for training and inference flops and holding the training dataset size fixed. The ability to enlarge patch size without restriction is a distinctive advantage of patch-centric models, liberating them from the efficiency constraints inherent to fixed-vocabulary token models, as outlined in Section~\ref{section:bpe-patch}. Using longer patches reduces computational demands, thereby allowing those resources to be redirected to expanding the global latent transformer, which is executed at a lower frequency as a result.

We perform a scaling analysis under fixed inference conditions to examine whether increasing model size while decreasing the number of steps on larger input patches leads to improved performance compared to smaller models using more steps. Beginning with baseline models containing 400 million and 3.6 billion parameters and utilizing the Llama 2 tokenizer, we identify flop-matched counterparts that use the Llama 3 tokenizer as well as {\textsc BLT}-Entropy architectures employing mean patch sizes of 6 and 8 bytes on the training datamix (refer to~\autoref{tab:fixed-inf-params} for comprehensive model configurations). For the models with patch size 8, we employ 3 encoder layers rather than just 1. We train each configuration across a spectrum of training flop allocations.

\autoref{fig:fixed_inference_scaling} illustrates that {\textsc BLT} models exhibit superior scaling behavior compared to tokenization-based systems across both inference FLOP categories. For both classes, BPE models initially yield higher performance given limited training resources, but they are soon outperformed by {\textsc BLT} models, even just beyond the compute-optimal range. This suggests that allocating increased resources during one-off pretraining may result in a model that delivers stronger performance at a predetermined inference cost. The 8B model class, exemplified by Llama 3.1, serves as a clear illustration: such models are trained on datasets that are two orders of magnitude larger than the compute-optimal data quantity for their scale.

The threshold at which {\textsc BLT} outperforms token-based approaches has shifted marginally nearer to the compute-optimal region as one moves to models with higher flop budgets (from roughly 3x to 2.5x the optimal compute allocation). In the case of models with a patch size of 8, the scaling trend becomes steeper within this higher flop class, enabling these models to surpass alternatives at earlier stages. As noted in Section~\ref{section:parameter-matched-scaling}, increasing patch size tends to result in performance more comparable to BPE models, particularly as model size grows. This phenomenon can partially be explained by the reduced proportion of overall computation consumed by the byte-level Encoder and Decoder, which do not scale as quickly as the Latent Transformer component. Notably, when expanding the model’s parameters by a factor of 20, from 400M to 8B, only an approximate doubling in the local parameters of {\textsc BLT} occurs. This detail is critical because increasing the patch size impacts only the Latent Transformer’s computational cost, leaving the flops of byte-level modules mostly unchanged. As a direct consequence, {\textsc BLT}-Entropy with ps=8 increased from constituting 1.6x to 1.7x the size of the Llama 2 model as model scale increased.

In conclusion, our investigation into patch-length scaling reveals that the {\textsc BLT} patch-based architecture benefits from concurrent increases in both patch size and model capacity, resulting in superior scaling performance. Notably, these positive scaling effects continue and are further enhanced as the model size grows.

\section{Byte Modeling Improves Robustness} Additionally, we assess the robustness of {\textsc BLT} relative to token-based models that do not access byte-level information directly, and introduce a method for augmenting pretrained token-based models to operate at the byte level.
\label{sec:robustness}

\subsection{Character-Level Tasks}

An initial impetus for developing byte-level models stemmed from their ability to handle noise at the byte level in input data, as well as their inherent sensitivity to the internal structure of tokens—a capability where traditional tokenizer-based models often fall short. To systematically assess these aspects, we conduct further experiments using benchmarks specifically designed to test both resilience to noisy input and fine-grained character-level comprehension. These include tasks covering English and multilingual scripts, along with categories such as digits and phonemes. The outcomes of these assessments are summarized in Table~\ref{tab:char_tasks}.

\paragraph{Noisy Data} To evaluate the resilience of models to noisy input, we generate perturbed versions of the benchmark classification datasets introduced in Section~\ref{section:evals}, facilitating a direct comparison between the robustness of tokenizer-based approaches and {\textsc BLT}. We introduce text variability using five separate character-level perturbation methods: (a) \textit{AntSpeak}: The entire text is rendered in uppercase, with characters separated by spaces. (b) \textit{Drop}: Ten percent of characters within the text are randomly deleted. (c) \textit{RandomCase}: Each character in the text is randomly assigned to be either uppercase or lowercase, resulting in a 50\% distribution for each. (d) \textit{Repeat}: Twenty percent of the characters are repeated up to four times each. (e) \textit{UpperCase}: Every character is changed to uppercase. For assessment, each noising strategy is individually applied to either the prompt, the completion, or both, with the final results averaged across these tasks. Table~\ref{tab:char_tasks} presents the performance on the noised HellaSwag~\citep{Eval_hellaswag} benchmark, demonstrating that {\textsc BLT} consistently surpasses tokenizer-based models in robustness, achieving on average an 8-point lead over models trained on identical data, and even exceeding the performance of Llama 3.1 trained on a considerably larger corpus.

\paragraph{Phonology - Grapheme-to-Phoneme (G2P)} We evaluate how effectively {\textsc BLT} translates sequences of graphemes (individual letters forming a word) into their corresponding phonemic representations (the spoken sounds of the word). Table \ref{tab:char_tasks} displays the outcomes for the G2P task under a 5-shot experimental setup, utilizing data from Phonology Bench~\citep{suvarna2024phonologybench}. Our results indicate that {\textsc BLT} achieves superior performance compared to the Llama 3 1T tokenizer-based baseline.

\paragraph{CUTE}
To evaluate character-level reasoning, we apply {\textsc BLT} to the CUTE benchmark~\citep{edman2024cute}, which consists of a variety of tasks divided into three main groups: compositionality understanding, orthographic similarity recognition, and sequence manipulation proficiency. This suite of tasks represents a notable challenge for typical tokenizer-based models, which tend to retain only limited practical knowledge of token spellings, impeding their capacity to manipulate textual sequences accurately. As presented in Table~\ref{tab:char_tasks}, {\textsc BLT}-Entropy surpasses both BPE Llama 3 variants by over 25 points on this suite of tasks. Notably, the model attains near-perfect accuracy (99.9\%) in character manipulation and spelling challenges. These substantial gains are achieved even though {\textsc BLT} was trained on only 1/16th the data volume compared to Llama 3.1, suggesting that BPE-based models inherently face difficulties acquiring robust character-level representations. Figure~\ref{fig:cute_examples} depicts several cases in which the Llama 3 tokenizer struggles, whereas our {\textsc BLT} system is notably effective. While BPE outperforms on the word deletion and insertion tasks, byte-level architectures may find these word-level manipulations more challenging; still, the performance gap remains modest, and incrementally building from characters to words could prove more tractable than the reverse direction. All experiments use a unified evaluation configuration and the original Huggingface prompts. It should be noted that BPE architectures might yield improved performance with further prompt engineering.

\paragraph{Low Resource Machine Translation} We assess the performance of {\textsc BLT} for translation tasks involving both directions—into and out of—six widely spoken language families and a set of twenty-one low-resource languages with diverse scripts, utilizing the FLORES-101 benchmark~\citep{goyal2022flores}. SentencePiece BLEU results are summarized in Table~\ref{tab:flores}. The findings reveal that {\textsc BLT} surpasses the version using the Llama 3 tokenizer, providing an overall increase of 2 BLEU points for translation into English and an improvement of 0.5 BLEU points for translation from English. In language pairs with extensive resources, {\textsc BLT}'s performance is on par with, or marginally superior to, Llama 3. Notably, for a wide range of language pairs among lower-resource families, {\textsc BLT} consistently achieves better results than Llama 3, highlighting the strengths of byte-level modeling when handling rare or long-tail byte sequences.

\subsection{Training {\textsc BLT} Using Llama 3 as Initialization}
\label{sec:distillation}
We investigate a methodology whereby {\textsc BLT} models benefit from prior knowledge embedded in existing pre-trained, tokenizer-based models to facilitate improved and accelerated convergence during training. This is implemented by initializing the global transformer parameters of {\textsc BLT} using those from a pre-trained Llama 3.1 checkpoint. In this procedure, we fine-tune the global transformer parameters at a learning rate that is one-tenth of that applied to the local encoder and decoder, maintaining the same number of optimization steps determined as optimal for Llama 3. The empirical results in Table~\ref{tab:distill} offer a performance comparison with a baseline version of {\textsc BLT}. The findings demonstrate that initializing {\textsc BLT} from Llama 3.1 leads to marked improvements, surpassing both the original Llama 3 and baseline {\textsc BLT} models, all trained within comparable computational constraints. Furthermore, in contrast to our {\textsc BLT}-Entropy model—which, as shown in Table~\ref{tab:evals}, was trained on a substantially larger set of 1T tokens—the Llama 3.1-initialized {\textsc BLT} attains higher scores on the MMLU benchmark. This indicates that leveraging pre-trained models for initialization can be a powerful strategy to markedly cut down on total training compute without sacrificing performance.

This approach may be interpreted as adapting tokenizer-based models to operate without tokenizers, essentially transforming a pre-trained LLaMA 3.1 model into a {\textsc BLT} framework. For a thorough evaluation, we present the original LLaMA 3.1—trained on 15T tokens—in Table \ref{tab:distill} and compare its results with those of the {\textsc BLT} version derived from LLaMA 3. Our findings indicate minor performance degradation on MMLU and HumanEval tasks, while observing a more pronounced decrease across other benchmarks. These results highlight the necessity for additional research to make optimal use of pre-trained models and enhance their performance, particularly by refining the selection of data mixtures and other critical hyperparameters.

\section{Ablations and Discussion}
\label{section:ablations}
This section presents ablation studies that substantiate our architectural selections for {\textsc BLT}, alongside analysis of the patching method and key hyper-parameters applied to the {\textsc BLT} 8B-parameter model trained on the {\textsc BLT}-1T dataset.

\paragraph{Entropy Model Hyper-parameters} To analyze how entropy model capacity and context window size affect scaling outcomes, we experiment with byte-level entropy transformer models ranging from 1m to 100m parameters and vary context windows from 64 up to 512 tokens. We display training flop scaling law plots of bits-per-byte (bpb) for our $400m$ and $1b$ {\textsc BLT} models, each trained on the Llama-2 dataset; results can be found in \autoref{fig:entropy_ablation}. Our findings indicate that increasing either the size of the entropy model or the length of its context generally leads to better scaling behavior, though the improvements level off beyond 50m parameters.

\paragraph{Types of Patching}

We conduct an ablation study on the four patching methodologies outlined in Section~\ref{section:patching}: 1) Strided Patching with strides of 4 and 6, 2) Whitespace-based Patching, 3) Patching utilizing the BPE Tokenizer associated with the Llama 3 tokenizer, and 4) Entropy-driven patching employing a compact byte-level LLM.

Although dynamic patching shortens the actual sequence lengths, we ensure that the context length remains consistent across all patching methods by regulating the sequence length. This allows all models, on average, to process an equal number of bytes per sequence during both training and inference, thereby eliminating potential confounds related to different context sizes. The outcomes of these ablation experiments are shown in Figure~\ref{fig:scaling-compute-opt}. Notably, every patching approach other than static patching delivers superior performance, with space patching closely rivaling dynamic entropy based patching.

\autoref{tab:evals_ablation} displays the results of benchmark evaluations for {\textsc BLT} models, providing a comparison between tokenizer-based approaches, space patching, and entropy-based patching. These models are trained on the Llama 2 dataset, each for a number of steps deemed optimal~\citep{dubey2024llama}. While space patching constitutes a more straightforward method that circumvents the need for an on-the-fly entropy model during training, our analysis reveals that the improvements observed with entropy-based patching in scaling experiments~(Section~\ref{section:scaling}) persist and are also reflected in downstream benchmark performance.\footnote{Results for space patching are taken from prior experiments conducted without cross-attention; however, similar patterns have been noted when cross-attention is included.}

\paragraph{Cross-Attention} 
\autoref{tab:crossattn_ablations_1b} presents an ablation study on the effects of incorporating cross-attention at different stages within the encoder and decoder of {\textsc BLT}. Specifically, for cross-attention in the encoder, we compare three distinct strategies for initializing the queries: (1) a uniform learned embedding applied across all global states, (2) an embedding generated via hashing the patch bytes, and (3) an aggregated, or pooled, version of the encoder's hidden representation of the patch bytes at the particular encoder layer.

Our results indicate that implementing cross-attention within the \textit{decoder} yields the highest efficacy. Incorporating cross-attention into the encoder offers modest benefits, but these gains are limited to scenarios where queries are initialized via pooling. Furthermore, cross-attention demonstrates notable advantages when applied to Common-Crawl data, with the effect being especially pronounced for models employing larger patch sizes.

\paragraph{n-gram Hash Embeddings} 

We explore various configurations of n-gram hash embedding vocabularies—specifically, 0, 100K, 200K, and 400K sizes—and summarize the outcomes in Table~\ref{tab:ngram_ablations_1b}. Our analysis shows that introducing hash embeddings yields benefits across all tested domains, with the largest impact observed on Wikipedia and Github (a 0.04 bpb improvement, compared to a 0.01 bpb improvement after 15k steps at the 8B scale). At 8B scale, decreasing the number of hashes from 500K to 300K produces a marginal effect on performance (about a 0.001 bpb difference at 15k steps). These findings suggest that utilizing hashes is crucial for aligning the performance of {\textsc BLT} with tokenizer-based approaches; however, increases beyond 300K hashes yield limited additional benefit. Furthermore, the improvements conferred by hash embeddings appear to complement those from cross-attention, as both approaches enhance results on distinct datasets.

\paragraph{Local Model Hyperparameters} Table \ref{tab:local_layers} presents an ablation study exploring different choices for the layer count in both the local encoder and decoder. Results indicate that, in combination with hash n-gram embeddings, {\textsc BLT} achieves strong performance with a highly minimal encoder—using only a single layer—while benefiting from a more substantial, deeper decoder.

\section{Related Work}
\label{section:related_work}

\textbf{Character-Level RNNs:} The task of Character Language Modeling has maintained its popularity since the inception of early neural architectures~\citep{sutskever2011generating,mikolov2012subword,graves2013generating} due to their inherent ability to naturally handle out-of-vocabulary terms, eliminating the need for traditional back-off techniques. In the work by \cite{kim2016character}, a model is introduced that relies solely on input characters, employing a combination of convolutional layers and highway networks as input to LSTM-driven RNNs. This architecture achieved comparable results to then state-of-the-art RNN-based language models in English, and demonstrated superior effectiveness on languages with complex morphology—an additional reason for the sustained interest in character-level LLMs. Further supporting this trend, \cite{kenter2018byte} employed byte-level LSTM models for machine comprehension and surpassed word-based models in performance on morphologically rich languages such as Turkish and Russian. Along similar lines, \cite{zhang2015character} leveraged character-based convolutional models in classification tasks, occasionally obtaining better results than those built on word-level representations. Hierarchical LSTM structures, as proposed by \citet{chung2019hierarchical}, integrate boundary detectors at various levels to discern latent textual hierarchies, thereby improving outcomes in character-level language modeling. The ByteNet model from \cite{kalchbrenner2016neural} instead utilizes CNN layers operating at the character level, offering an alternative to attention mechanisms in machine translation tasks.

\textbf{Character-Level Transformers:} Transformer architectures leveraging attention mechanisms~\citep{vaswani2017attention} combined with subword tokenization techniques~\citep{sennrich-etal-2016-neural} have resulted in substantial advancements in neural language modeling and the attainment of state-of-the-art performance on standard benchmarks. Nevertheless, employing word or sub-word tokens introduces a built-in inductive bias concerning the level at which linguistic abstraction is modeled. Seeking to integrate the strengths of transformers with early achievements in character-level language modeling, \citet{al2019character} employ exceptionally deep transformer architectures. By incorporating auxiliary loss functions, they succeed in training transformer models that achieve better results than former LSTM-based character language models. Despite this progress, a notable performance disparity relative to word-level LLMs persisted. Similarly, GPT-2~\citep{radfordgpt2} found that on extensive datasets such as the 1 Billion Word benchmark, models operating at the byte level fell short of matching the effectiveness of word-level language models.

Previous work by \cite{Choe2019BridgingTG} showed that transformer-based byte-level LLMs can surpass subword-level counterparts with similar parameter counts, but at the cost of significantly higher computational requirements and prolonged training times. In a similar vein, \cite{el2020characterbert} introduced CharFormer, a BERT-based model that constructs word-level representations through convolutional operations on character embeddings; although they report notable performance gains within the medical domain, this comes alongside a substantial increase in computational overhead. The CANINE architecture introduced by \cite{clark2022canine}, with 150M parameters, directly processes character sequences through an encoder-only approach. Like our global model, CANINE utilizes a deep transformer stack in addition to a local transformer and strided convolutions to perform input downsampling, leading to superior performance over comparable token-level encoder-only models, such as mBERT, particularly on multilingual downstream benchmarks. ByT5~\citep{xue2022byt5} investigated byte-level encoder-decoder models that omit patching operations altogether; while their system demonstrated heightened resilience to input noise and matched tokenization-based models using only a quarter of the training data, the absence of patching required performing costly attention over every individual byte, resulting in substantial computational demands. Transitioning from subword units to bytes increases sequence lengths, thereby complicating efficient scaling for byte-level language models. More recently, leveraging the Mamba Architecture~\citep{gu2023mamba}—which supports a fixed-size memory over extended contexts—\cite{wang2024mambabyte} successfully trained a byte-level Mamba model, likewise forgoing patching. In a computation-constrained setup at the 350M parameter level, this architecture surpassed the performance of byte-level transformers in bits-per-byte across several datasets.

\textbf{Patching-based approaches:} Utilizing patching methods has proven effective in reducing the substantial computational demands typically associated with byte-level LLMs, without necessarily sacrificing performance. Numerous studies have presented promising results at relatively modest model scales and training data volumes. For example, \cite{nawrot-etal-2022-hierarchical} explored static approaches to patching by leveraging downsampling and upsampling, introducing the hourglass transformer, which surpassed competing byte-level models at the 150M scale. Building on this, \cite{nawrot-etal-2023-efficient} introduced dynamic patching mechanisms. Their enhancements included an end-to-end learnable boundary-predictor, a version trained with the guidance of specific tokenizers, as well as an entropy-driven patching system akin to {\textsc BLT}. Their findings indicate these techniques surpassed standard transformer architectures in language modeling benchmarks, using models with 40M parameters on datasets of approximately 400M tokens. Separately, \cite{lester2024training} explored sequence compression with arithmetic coding, achieving higher compression ratios than BPE, and applied the equal-info windows technique. This approach yielded superior results compared to byte-level baselines in language modeling tasks, although it did not reach the performance of subword-based baselines.

Our research is principally influenced by, and most comparable to, MegaByte~\citep{yu2023megabyte}, a decoder-only causal LLM. MegaByte leverages a predetermined patching and representation concatenation technique to transform bytes into patches; subsequently, a localized decoder-side model reconstructs bytes from these patches. The authors show that at a 1B parameter scale, MegaByte achieves performance on par with models that rely on tokenizers, tested over approximately 400B bytes of data. In our evaluations, we consistently include MegaByte and observe that this static patching methodology underperforms relative to the best compute-efficient, tokenizer-based models when flops are held constant. We show that {\textsc BLT} effectively addresses this performance disparity. Similarly, \cite{slagle2024spacebyte} report that MegaByte falls short and propose enhancements such as adapting the patching mechanism to whitespace and other space-character bytes, along with introducing a local encoder. Their findings indicate certain gains over tokenizer-based transformers within specific domains—like Github and arXiv—at the 1B parameter scale in a compute-constrained context. We extend these investigations with experiments using their model, ultimately demonstrating that further advances in architecture are necessary for byte-level models to successfully reach the performance exhibited by leading token-based models like Llama 3.

\section{Limitations and Future Work}

For this study, our approach to selecting model architectures involves training for the same optimal step count identified for Llama~3~\citep{dubey2024llama}. Nevertheless, since these scaling laws are specifically derived for BPE-level transformers, their direct application might not yield ideal data-to-parameter ratios when applied to the {\textsc BLT} setting. A comprehensive derivation of scaling laws tailored to {\textsc BLT} is deferred to future research; such an undertaking could potentially reveal even more advantageous scaling properties for this model family. Moreover, the current experimental configurations predominantly target models with up to 1B parameters, suggesting that optimal architectural configurations may shift as the model size increases to 8B parameters or more, which might further enhance performance at larger scales.

Current transformer frameworks and repositories are predominantly optimized for transformer architectures that rely on tokenization. Although our experimental setup includes theoretical flop-equivalent comparisons and leverages efficient solutions like FlexAttention to accommodate layers that depart from the standard transformer structure, our implementations may still lag behind tokenizer-centric models in terms of wall-clock performance. Consequently, additional optimizations could further enhance their computational efficiency.

Although {\textsc BLT} relies on a separately trained entropy model for patching, exploring the possibility of jointly learning the patching model within an end-to-end framework represents a promising avenue for future research. In Section \ref{sec:distillation}, we provide preliminary results suggesting that it is feasible to “byte-ify” tokenizer-based models—such as Llama 3, which have been trained on over 10T tokens—by initializing the global transformer with their pretrained weights and subsequently freezing it. Advancing research along these lines could yield strategies that both preserve the advantages of bytefying and achieve higher performance than these tokenizer-based models, all without the need for retraining from scratch.

\section{Conclusion}
\label{section:conclusion}

In this work, we introduce the Byte Latent Transformer (\textbf{BLT}), an architectural innovation that challenges the traditional reliance on fixed-vocabulary tokenization in large language models. {\textsc BLT} employs a dynamic, trainable approach to aggregating bytes into patches, thereby enabling the adaptive allocation of computational effort according to the underlying data's complexity. This strategy yields notable gains in both computational efficiency and model robustness. Through comprehensive scaling experiments, we show that {\textsc BLT} can achieve parity with tokenization-based architectures such as Llama 3 at scales up to 8B parameters and 4T bytes, while potentially reducing inference computational costs by as much as 50\% in exchange for slight declines in standard evaluation scores. Additionally, {\textsc BLT} paves the way for a novel scaling approach, permitting simultaneous expansion of model and patch dimensions without exceeding a predefined inference resource constraint—a benefit particularly suited for practical computing scenarios. By processing raw byte sequences directly, {\textsc BLT} not only enhances resilience to atypical or corrupted data but also fosters richer representations of sub-word patterns. Collectively, these findings underscore {\textsc BLT}'s promise as an efficient, resilient, and extensible alternative to established tokenization-centered language modeling paradigms.

\end{document}